% A made-up paper, in two versions, to show what arxiv-diff reports. Nothing here is a real result.
\documentclass{article}
\usepackage{booktabs}
\newcommand{\method}{TinyFormer}
\title{\method: Small Attention Models for Toy Translation}
\author{Ada Lovelace \\ Analytical Engines Ltd \And Charles Babbage \\ Difference Lab}
\begin{document}
\maketitle
\begin{abstract}
We present \method, a small attention model for translating toy languages.
It reaches 31.2 BLEU on ToyWMT, 2.1 points above the best recurrent model.
\end{abstract}

\section{Introduction}
Recurrent models translate toy languages well but slowly~\cite{elman1990}.
We ask whether attention alone suffices at small scale.

\section{Method}
\method{} stacks four attention layers of width 256.
Each layer has eight heads.

\section{Experiments}
We train on ToyWMT for 20 epochs.
Table~\ref{tab:main} shows that \method{} reaches 31.2 BLEU.

\begin{table}[t]
\centering
\caption{BLEU on the ToyWMT test set.}\label{tab:main}
\begin{tabular}{lcc}
\toprule
Model & EN-XX & XX-EN \\
\midrule
GRU & 29.1 & 27.8 \\
LSTM & 28.7 & 28.0 \\
\method{} & \textbf{31.2} & \textbf{29.4} \\
\bottomrule
\end{tabular}
\end{table}

\section{Conclusion}
Attention alone is enough for toy translation.
We will release the code soon.

\bibliographystyle{plain}
\bibliography{refs}
\end{document}
